\documentclass[aspectratio=169,10pt,english]{beamer}

\input{../../../_preambulo_beamer.tex}
\input{../../../_comandos_beamer.tex}

\selectlanguage{english}

\title{MA1001B - Statistical Modeling for Decision Making}
\subtitle{Lesson 30: Sample Size for Margin of Error and Budget}
\author{}
\institute{Tecnol\'ogico de Monterrey}
\date{}

\begin{document}

\begin{frame}[plain]
  \vspace{1.2cm}
  {\Large\bfseries MA1001B - Statistical Modeling for Decision Making\par}
  \vspace{0.7cm}
  {\large Lesson 30: Sample Size for Margin of Error and Budget\par}
  \vspace{0.7cm}
  {\color{TecRojo}\rule{\textwidth}{1.2pt}}
  \vspace{0.8cm}

  MA1001B Faculty\\[0.3cm]
  Term\\[0.3cm]
  Tecnol\'ogico de Monterrey
\end{frame}

\begin{frame}{The Planning Question}
  \begin{alertblock}{Guiding question}
    How large must $n$ be so that a 95\% interval hits a chosen margin of
    error, and can the budget fund that $n$?
  \end{alertblock}

  \vspace{0.6em}
  By the end of this lesson, you can:
  \begin{itemize}
    \item compute $n=(z^{*}s/E)^{2}$ for a mean and round up;
    \item compute $n=z^{*2}p(1-p)/E^{2}$ for a proportion;
    \item explain why $p=0.5$ is conservative;
    \item compare required $n$ with the $n$ a budget can buy.
  \end{itemize}

  \vfill
  \scriptsize All planning values used today are synthetic. No real company data are used.
\end{frame}

\begin{frame}{Choose $E$ Before Collecting Data}
  \small
  \begin{center}
    \begin{tabular}{lr}
      \toprule
      \textbf{Mean-handling plan} & \textbf{Value} \\
      \midrule
      Textbook $z^{*}$ & 1.96 \\
      Planning $s$ & 10 minutes \\
      Target margin $E$ & 2 minutes \\
      \bottomrule
    \end{tabular}
  \end{center}

  \vspace{0.5em}
  \begin{exampleblock}{Mean sample-size formula}
    \[
      n=\left(\frac{z^{*}s}{E}\right)^{2}=\left(\frac{1.96\times 10}{2}\right)^{2}=96.040000.
    \]
    Round up to the next whole observation: $n=97$.
  \end{exampleblock}
\end{frame}

\begin{frame}[fragile]{Step 1: Sample Size for a Mean}
  \begin{columns}[T]
    \begin{column}{0.42\textwidth}
      \small
      \begin{lstlisting}
raw = (z_star * s / E) ** 2
n = int(np.ceil(raw))
      \end{lstlisting}

      \begin{block}{Verified output}
        Raw $n=96.040000$\\
        Rounded-up $n=97$
      \end{block}
    \end{column}
    \begin{column}{0.58\textwidth}
      \centering
      \includegraphics[width=\textwidth]{../figures/sample_size_01_mean_margin.png}
      \vspace{0.2em}
      \scriptsize Required $n$ versus $E$ from Step 1
    \end{column}
  \end{columns}
\end{frame}

\begin{frame}{Conservative $p=0.5$ for a Proportion}
  \small
  Target margin $E=0.03$ for a first-contact rate:
  \[
    n=\frac{z^{*2}p(1-p)}{E^{2}}.
  \]
  At $p=0.5$, $p(1-p)=0.25$ is maximized:
  \[
    n=\frac{1.96^{2}\times 0.25}{0.03^{2}}=1067.111111 \;\to\; 1068.
  \]
  Using the Lesson 29 guess $p=0.18$ yields only $n=631$, which is 437
  observations too optimistic if $p$ is nearer $0.5$.
\end{frame}

\begin{frame}[fragile]{Step 2: Conservative Proportion Sample Size}
  \begin{columns}[T]
    \begin{column}{0.40\textwidth}
      \small
      \begin{lstlisting}
raw = (
    z_star**2 * p * (1-p)
    / E**2
)
      \end{lstlisting}

      \begin{block}{Verified output}
        $p=0.5$ gives $n=1068$\\
        $p=0.18$ gives $n=631$\\
        Extra $=437$
      \end{block}
    \end{column}
    \begin{column}{0.60\textwidth}
      \centering
      \includegraphics[width=\textwidth]{../figures/sample_size_02_proportion_conservative.png}
      \vspace{0.2em}
      \scriptsize $n$ versus planning $p$ from Step 2
    \end{column}
  \end{columns}
\end{frame}

\begin{frame}{Step 3: Desired $E=0.01$ Can Exceed the Budget}
  \begin{columns}[T]
    \begin{column}{0.42\textwidth}
      \small
      Conservative $n$ for $E=0.01$:
      \[
        n=9604.
      \]
      Budget $\$20{,}000$ at $\$12$ each
      funds only $1666$ observations.

      \vspace{0.4em}
      Required cost: $\$115{,}248$.

      \vspace{0.5em}
      \textbf{Limit:} shrinking $E$ inflates $n$ with $1/E^{2}$. A desired
      margin is not a fundable study.
    \end{column}
    \begin{column}{0.58\textwidth}
      \centering
      \includegraphics[width=\textwidth]{../figures/sample_size_03_budget_vs_precision.png}
      \vspace{0.2em}
      \scriptsize Required versus affordable $n$ from Step 3
    \end{column}
  \end{columns}
\end{frame}

\begin{frame}{Concept Check}
  \begin{enumerate}
    \item Compute $(1.96\times 10/2)^{2}$ and round up.
    \item Why is $p=0.5$ conservative for a proportion?
    \item Why does $n$ explode when $E$ shrinks from $0.03$ to $0.01$?
    \item Can a $\$20{,}000$ budget at $\$12$ per observation fund $n=9604$?
  \end{enumerate}

  \vspace{0.6em}
  \begin{block}{Connection to Lesson 31}
    The session can continue directly into a chi-square interval for a
    population variance.
  \end{block}

  \vfill
  \scriptsize Source: OpenStax, \textit{Introductory Business Statistics 2e},
  Chapter 8, Section 8.4. CC BY-NC-SA.
\end{frame}

\appendix



\end{document}
